\documentclass[letterpaper]{article}

% IROS 2026 Format - Based on IEEE Conference Proceedings
\usepackage{cite}
\usepackage{amsmath,amssymb,amsfonts}
\usepackage{algorithm}
\usepackage{algorithmic}
\usepackage{graphicx}
\usepackage{textcomp}
\usepackage{xcolor}
\usepackage{booktabs}
\usepackage{multirow}
\usepackage{array}
\usepackage{caption}
\usepackage{geometry}
\usepackage{hyperref}

% Page margins for IROS format
\geometry{
    top=1in,
    bottom=1in,
    left=0.75in,
    right=0.75in
}

\begin{document}

\title{Gaze2Guide: Memory-Augmented Gaze Reasoning for Exhibit-to-Exhibit Decision and Dwell-Time Prediction in Real Exhibition Spaces}
\author{Xijing Wang, Lei He, and Nianxiong Liu}
\date{}

\maketitle

\begin{abstract}
Museum guide and service robots need human goal and timing priors to decide where to guide and when to intervene without being intrusive. We propose IROS\_gaze, a memory-augmented gaze reasoning framework that predicts (i) the visitor's next exhibit (topological goal), (ii) dwell time and an ordinal attention level, and (iii) a K-step rollout of future viewing trajectories. The system integrates exhibit semantics, an exhibit-level neighbor graph, and real eye-tracking observations through a modular pipeline: ContextBuilder, a dual-layer Memory (short-term window N=5 and long-term statistics), an LLMReasoner, and an iterative SequencePredictor (K=5). We construct instruction-style training data from real eye-tracking collected in two venues (Tsinghua Art Museum and an Osaka exhibition), resulting in 8,159 high-quality labeled segments after cleaning from 9,158 raw records. Across three tasks (scan-path planning, next-step prediction, and attribution), we report improvements over strong baselines under both random and cross-subject splits, and demonstrate a robot-facing topological guide simulation where the learned priors reduce regret and improve intervention timing.
\end{abstract}

\section{INTRODUCTION}

Robots operating in public indoor environments (e.g., museums and exhibitions) must anticipate where a visitor will go next and how long they will stay, to plan socially appropriate guidance, explanation, and collision-free motion. A guide robot that can predict visitor intent and timing gains the ability to: (1) position itself proactively at the next viewing location, (2) time its explanations to coincide with natural attention peaks, and (3) avoid interrupting high-engagement moments. While gaze is a rich signal of human attention, prior work often uses it only for visualization (heatmaps) or isolated intent cues, without producing robot-consumable goal and timing priors constrained by the environment's topology.

\subsection{Motivation: Human-Aware Robot Guidance}

Consider a museum guide robot that must decide when to approach a visitor to provide information. If the robot interrupts too early, it disrupts the visitor's engagement with the current exhibit. If it waits too long, the visitor may have already moved on. Similarly, for following behavior, the robot must predict which exhibit the visitor will approach next to position itself appropriately without blocking the viewing angle. These tasks require understanding both \textit{where} the visitor will go (spatial goal) and \textit{how long} they will stay (temporal engagement).

Existing approaches to human intention prediction often rely on trajectory extrapolation or spatial priors alone. However, in structured environments like museums, visitor behavior is heavily influenced by semantic factors—exhibit content, thematic relationships, and personal interests—that are not captured by geometric reasoning alone. Eye-tracking provides a rich signal of attention and cognitive processing, but prior work has not fully exploited it for robot-consumable prediction in topological spaces.

\subsection{Our Approach}

We treat gaze as a cognitive sensor for embodied spatial reasoning. Our key idea is to combine (i) an exhibit-level neighbor graph encoding spatial layout, (ii) short- and long-term memory of viewing history, and (iii) LLM-based constrained reasoning, to jointly predict next exhibit, dwell time, and attention level, and to roll out future viewing trajectories. This yields a plug-in prior module that can bias a guide agent's action selection and intervention timing.

Our framework, IROS\_gaze, operates at the exhibit level rather than pixel-level gaze prediction. This topological abstraction is both robot-relevant (robots navigate between exhibits, not pixels) and data-efficient (exhibit-level labels can be collected via simple hotkey annotations during eye-tracking sessions).

\subsection{Contributions}

\textbf{Our contributions are:}

\begin{enumerate}
    \item \textbf{A robot-ready gaze-driven prior:} A framework that learns from human gaze data to predict goal location, dwell time, and interaction readiness, producing outputs directly consumable by guide/service robot decision modules. Unlike pure visualization approaches, our outputs are structured probabilities over discrete goals with associated timing estimates.

    \item \textbf{A structured reasoning pipeline:} ContextBuilder fuses topology constraints with memory; dual-layer Memory (short-term N=5 window + long-term statistics) captures temporal patterns; LLMReasoner performs constrained CoT reasoning; SequencePredictor performs K=5 rollouts for planning. Each component is modular and can be adapted to other structured environments.

    \item \textbf{Robot-style evaluation:} Beyond prediction metrics, we validate in a topological guide simulation (Experiment-R) showing that our learned priors reduce policy regret and improve intervention timing compared to strong baselines. This simulation directly addresses robot-relevant questions: can the predictions help a guide robot make better decisions?
\end{enumerate}

\section{RELATED WORK}

\subsection{Gaze-Based Intention and Engagement Modeling}

Gaze has been widely used as a signal of human attention and intent in human-robot interaction. Early work by Shon et al. used gaze direction as a deictic reference for collaborative robots, enabling joint attention through gaze following. More recent approaches combine gaze with contextual cues such as head pose, facial expressions, and body language for engagement estimation. However, most prior work treats gaze as an instantaneous indicator, without modeling its cumulative relationship to visitation goals and timing in structured environments.

In the museum and gallery domain, gaze-based systems have focused primarily on attention visualization—generating heatmaps to show which artworks attracted the most viewing time. While valuable for curatorial analysis, these outputs do not directly support robot decision-making. Our work differs by predicting discrete future behaviors (next exhibit, dwell time) that can be consumed by robot planning algorithms.

\subsection{Egocentric Vision and Gaze Datasets}

Several egocentric vision datasets capture gaze during daily activities. Fathi et al. proposed learning to predict daily actions from eye gaze, establishing gaze as a precursor to action. Yan et al. studied gaze-based intent prediction in supermarket shopping scenarios. Kitani et al. introduced activity forecasting from first-person video.

However, these datasets primarily serve visualization or saliency modeling tasks, and few capture gaze in structured, topological environments like museums. Our work differs by constructing instruction-style training data for predictive reasoning under topological constraints, enabling LLM-based reasoning that combines semantic and spatial cues.

\subsection{Topological Navigation and Goal Prediction}

Topological representations have proven effective for robot navigation, particularly in large-scale environments where metric representations are costly to maintain. Walter et al. presented a learning-based framework for topological map building and planning. Goal prediction in indoor spaces has been explored using trajectory data and semantic cues. Ziebart et al. proposed intent-driven trajectory prediction using inverse optimal control.

We extend this work by incorporating gaze-driven attention as a strong signal for predicting goals in exhibit-level graphs. While traditional approaches rely on past trajectories alone, our method leverages real-time attention signals to predict goals before any physical movement occurs.

\subsection{Human-Aware and Socially-Aware Navigation}

Socially-aware navigation requires robots to predict human motion and avoid intrusion. Trautman and Krause proposed unsupervised learning of cost functions for navigation in crowds. Chen et al. introduced socially aware motion planning with deep reinforcement learning. Human intention prediction has been studied for assistance robots in collaborative settings.

Our work contributes to this area by providing a gaze-based prior for goal and timing prediction that can be integrated into social navigation planners. Unlike trajectory-based prediction, our approach works from static gaze observations, enabling prediction before the visitor moves.

\subsection{Language Models for Constrained Reasoning}

Recent work demonstrates LLMs' capability for spatial reasoning and planning under constraints. Xia et al. showed that chain-of-thought prompting enables effective planning with large language models. Liang et al. surveyed LLM capabilities for spatial reasoning tasks.

Our LLMReasoner leverages these capabilities with specialized prompts that enforce topological feasibility and interpret gaze patterns through semantic and temporal coherence. By fine-tuning on instruction-style data, we achieve better performance than zero-shot baselines while maintaining the reasoning capabilities of large language models.

\section{DATA AND PROBLEM SETUP}

\subsection{Multi-Modal Inputs}

Each training segment corresponds to an exhibit visit interval. Inputs include:

\begin{enumerate}
    \item \textbf{Exhibit Semantic Descriptions:} Textual descriptions of exhibits, including visual features, historical context, and thematic relationships. These are obtained from exhibit labels and catalog information, enriched by VLM-based visual feature extraction when exhibit photos are available.

    \item \textbf{Exhibit-Level Neighbor Graph:} A topological graph $\mathcal{G} = (\mathcal{V}, \mathcal{E})$ where $\mathcal{V}$ is the set of exhibit IDs and $\mathcal{E}$ encodes spatial adjacency. Each edge is annotated with a relation type: ``next-along-path'' (sequential viewing order), ``visible-from'' (line-of-sight), or ``same-zone'' (thematic grouping).

    \item \textbf{Eye-Tracking Signals:} Raw gaze data sampled at 30 Hz, including gaze points $(x, y)$ in egocentric image coordinates, fixation/saccade events from dispersion-based detection, validity flag (tracking confidence), and pupil/blink measurements. Gaze is aligned with egocentric video timestamps.
\end{enumerate}

\subsection{Dataset Construction}

We collected eye-tracking data in two venues with different characteristics:

\begin{table}[h]
\centering
\caption{Data Collection Summary}
\begin{tabular}{lcccc}
\toprule
Venue & Type & Duration & Participants & Exhibits \\
\midrule
Tsinghua Art Museum & Art exhibition & $\sim$45 min/session & 12 & 28 \\
Osaka Exhibition & Cultural exhibition & $\sim$30 min/session & 8 & 19 \\
\bottomrule
\end{tabular}
\end{table}

The Tsinghua Art Museum session featured traditional Chinese paintings and calligraphy, while the Osaka Exhibition focused on interactive cultural displays. This diversity helps evaluate cross-domain generalization.

\textbf{Exhibit Visit Segmentation:} We derive exhibit visit intervals using hotkey annotations aligned with an exhibit dictionary. During data collection, participants pressed designated hotkeys when starting to view a new exhibit, creating ground-truth boundaries. This method, while simple, produces clean labels aligned with human perception of exhibit transitions.

\textbf{Data Cleaning:} Raw eye-tracking records (30 Hz) yield 9,158 records. We apply the following filtering rules: (1) discard records with average validity $<$ 60\% (poor tracking quality), (2) remove visits $<$ 3 seconds (transitional glances rather than genuine viewing), (3) remove visits $>$ 300 seconds (likely tracking errors or extended breaks), (4) aggregate consecutive fixations within the same exhibit. After cleaning, we obtain \textbf{8,159 high-quality labeled segments}.

\textbf{Attention Level Discretization:} We discretize dwell time into an ordinal 5-level attention scale using an exponential decay mapping. This discretization serves two purposes: (1) it makes the prediction task more tractable by reducing the output space, and (2) it aligns with robot-relevant categories (e.g., ``high engagement'' means don't interrupt).

\begin{table}[h]
\centering
\caption{Attention Level Definition (Exponential Decay)}
\begin{tabular}{clc}
\toprule
Level & Duration Range & Interpretation \\
\midrule
A & 90--300s & Deep engagement, careful examination \\
B & 45--90s & Sustained attention, normal viewing \\
C & 20--45s & Moderate interest, browsing \\
D & 8--20s & Brief interest, scanning \\
E & 3--8s & Glancing, minimal attention \\
\bottomrule
\end{tabular}
\label{tab:attention}
\end{table}

The exponential decay reflects the natural distribution of viewing times: many brief glances, fewer sustained viewings, and rare extended examinations.

\subsection{Task Formulation}

We define three complementary tasks that cover different aspects of visitor behavior understanding:

\begin{itemize}
    \item \textbf{T1 (Scan-Path Planning):} Given a set of visible exhibits with semantic features, predict an ordered viewing sequence with attention levels. This mimics the robot's task of understanding a visitor's intended viewing path upon entering a new zone.

    \item \textbf{T2 (Next-Step Prediction):} Given current exhibit, gaze history (including dwell times and attention levels), and spatial context (reachable neighbors), predict the next exhibit and attention level. This is the core prediction for robot following/guiding.

    \item \textbf{T3 (Attribution):} Generate an explanation for why an exhibit receives a given attention level based on visual features, visitor history, and context. This provides interpretability for robot decisions and enables natural language interaction.
\end{itemize}

\begin{table}[h]
\centering
\caption{Dataset Statistics}
\begin{tabular}{lr}
\toprule
Statistic & Value \\
\midrule
Venues & Tsinghua Art Museum; Osaka Exhibition \\
Sampling Rate & 30 Hz \\
Raw / Clean Segments & 9,158 / 8,159 \\
Task Counts & T1: 3,038; T2: 3,038; T3: 2,083 \\
Unique Exhibits & 47 (TH: 28, OS: 19) \\
Avg. Sequence Length & 6.2 exhibits per session \\
Attention Distribution & A: 12\%, B: 28\%, C: 25\%, D: 22\%, E: 13\% \\
Avg. Transitions per Exhibit & 3.7 (outgoing) \\
\bottomrule
\end{tabular}
\label{tab:dataset}
\end{table}

\subsection{Evaluation Splits}

We report results on two evaluation protocols:

\textbf{Random Split:} 9:1 train/validation split (7,343 train / 816 validation). This evaluates the model's ability to generalize to new sequences from the same population.

\textbf{Leave-Subject-Out:} Train on N-1 participants, test on held-out participant. This evaluates generalization to new visitors with potentially different viewing patterns and preferences.

\section{METHOD}

\subsection{System Overview}

IROS\_gaze consists of four layers (Fig.~\ref{fig:architecture}): (1) Data Input Layer processes raw gaze and exhibit data; (2) Memory Management Layer maintains temporal patterns; (3) Reasoning \& Prediction Layer performs constrained reasoning and multi-step lookahead; (4) Robot Interface Layer outputs structured priors for robot decision-making.

\begin{figure}[h]
\centering
% \includegraphics[width=0.48\textwidth]{figures/architecture.pdf}
\fbox{\parbox{0.8\textwidth}{\centering \vspace{2cm} Figure 1: System Architecture \vspace{2cm}}}
\caption{System architecture. The IROS\_gaze pipeline takes gaze observations, exhibit semantics, and topology as input. ContextBuilder fuses these inputs; dual-layer memory maintains temporal patterns; LLMReasoner predicts next goal, dwell time, and attention level; SequencePredictor performs K=5 rollouts for planning.}
\label{fig:architecture}
\end{figure}

The modular design allows each component to be developed and evaluated independently. For example, Memory can be replaced with a more sophisticated implementation, or LLMReasoner can be swapped for a different base model.

\subsection{Context Builder}

ContextBuilder fuses viewing history with spatial topology to produce a structured context for reasoning. The context serves as input to LLMReasoner and must contain all relevant information for prediction.

Given current gaze record $g_t = (e_t, a_t, d_t)$ where $e_t$ is exhibit ID, $a_t$ is attention level, $d_t$ is duration, and history $\mathcal{H}_{t-1} = \{g_{t-N}, ..., g_{t-1}\}$, the context includes:

\textbf{Current Exhibit Information:} ID, name, semantic features (visual description, historical notes), current attention level, and estimated duration.

\textbf{Recent History:} The N=5 most recent gaze records, each with exhibit ID, attention level, and duration. This window captures immediate sequential patterns (e.g., ``scanning left to right'').

\textbf{Spatial Context:} Three types of topological information:
\begin{itemize}
    \item \textit{Previous path:} Exhibits visited before reaching the current exhibit, indicating the visitor's route through the exhibition.
    \item \textit{Next path:} Exhibits typically visited after the current exhibit, based on aggregate visitor flow.
    \item \textit{Reachable options:} Direct neighbors of the current exhibit, representing feasible next steps.
\end{itemize}

\textbf{Statistics:} Global summary including total gaze count, unique exhibits visited, visit counts per exhibit, and cumulative dwell times.

ContextBuilder formats this information as a structured prompt for LLMReasoner, ensuring spatial and temporal constraints are explicitly represented.

\subsection{Dual-Layer Memory}

\textbf{Short-Term Memory (STM):} Maintains a fixed-size sliding window of the most recent N=5 gaze records using a deque (double-ended queue) data structure. When a new observation arrives, it is added to the front, and the oldest record is removed if the window exceeds N=5.

STM captures immediate sequential patterns such as:
\begin{itemize}
    \item \textit{Scanning patterns:} Left-to-right or right-to-left progression through a wall of exhibits
    \item \textit{Re-examination:} Returning to a previous exhibit after viewing nearby alternatives
    \item \textit{Thematic coherence:} Consecutive visits to exhibits with related themes
\end{itemize}

The window size N=5 was chosen empirically: it captures enough history to establish patterns while remaining tractable for LLM reasoning.

\textbf{Long-Term Memory (LTM):} Maintains global statistics across the entire session:

\begin{itemize}
    \item \textit{Visit frequency:} $count(e_i)$ — how many times exhibit $e_i$ has been visited. Repeated visits indicate high interest.
    \item \textit{Cumulative dwell time:} $duration_{total}(e_i)$ — total time spent at exhibit $e_i$. Longer times suggest higher engagement.
    \item \textit{Attention distribution:} $p(a|e_i)$ — distribution of attention levels for exhibit $e_i$. Some exhibits consistently receive high attention.
    \item \textit{Transition frequencies:} $count(e_i \rightarrow e_j)$ — how often visitors move from $e_i$ to $e_j$. This encodes popular routes through the exhibition.
\end{itemize}

These statistics provide strong priors for both prediction (frequently visited exhibits are more likely to be revisited) and explanation (justifying attention levels based on historical patterns).

LTM is implemented as a set of dictionaries updated incrementally as new gaze records arrive. Queries to LTM return aggregated statistics in O(1) time.

\subsection{LLM Reasoner}

LLMReasoner performs constrained reasoning over the structured context to produce predictions. The reasoning is constrained in three ways:

\textbf{Topological Constraints:} The predicted next exhibit must be among the reachable neighbors from the current exhibit. This eliminates infeasible predictions (e.g., teleporting across the exhibition space).

\textbf{Semantic Coherence:} The prediction should respect thematic relationships. For example, if the visitor has been viewing landscape paintings, they are more likely to continue with landscapes rather than abruptly switching to abstract sculpture.

\textbf{Temporal Consistency:} The prediction should account for the visitor's engagement level. High attention at the current exhibit may suggest either (a) continued interest in similar exhibits or (b) fatigue and a desire for variety.

The prompt template explicitly encodes these constraints and asks the LLM to reason through them before producing a prediction. We use Chain-of-Thought prompting: the LLM first generates a reasoning trace explaining its prediction, then outputs the structured prediction.

\subsection{Sequence Predictor}

SequencePredictor performs K-step lookahead by iteratively applying the single-step predictor:

\begin{algorithm}[h]
\caption{K-Step Sequence Prediction}
\begin{algorithmic}[1]
\STATE trajectory $\leftarrow$ []
\STATE context $\leftarrow$ current\_context
\FOR{$k = 1$ to $K=5$}
    \STATE prediction $\leftarrow$ LLMReasoner.predict(context)
    \STATE trajectory.append(prediction)
    \STATE context $\leftarrow$ update(context, prediction)
    \IF{should\_terminate(context)} \textbf{break} \ENDIF
\ENDFOR
\RETURN trajectory
\end{algorithmic}
\end{algorithm}

The update step adds the predicted exhibit to the history, advances the ``current position,'' and refreshes the reachable neighbors. Termination conditions include: (a) reaching an exit zone, (b) predicting a loop (returning to a recently visited exhibit), or (c) confidence falling below a threshold.

This K=5 rollout enables the robot to plan ahead. For example, if the rollout predicts the visitor will view exhibits A $\rightarrow$ B $\rightarrow$ C in the next 3 steps, the robot can position itself near exhibit C to provide timely information, or it can choose an alternative route that doesn't interfere.

\subsection{Training}

We use instruction tuning with ShareGPT format. Each training example consists of a structured input (current exhibit, history, context) and a structured output (prediction with reasoning). The instruction format encourages the model to follow the specified prediction format.

We fine-tune Qwen-32B using LoRA (Low-Rank Adaptation) to make training computationally feasible:
\begin{itemize}
    \item Learning rate: 2e-4
    \item Batch size: 16 (with gradient accumulation over 4 steps for effective batch size 64)
    \item Epochs: 3 (early stopping based on validation loss)
    \item Max sequence length: 2048 tokens (sufficient for context + prediction)
    \item LoRA rank: 8 (balancing adaptability and efficiency)
\end{itemize}

Training converges in approximately 6 hours on a single A100 GPU.

\subsection{Robot-Facing Interface}

At each step $t$, the system outputs robot-consumable priors:

\begin{itemize}
    \item \texttt{goal\_distribution}: A probability distribution over reachable neighbors, $p(e_{t+1} | context)$. The robot can use this for probabilistic planning or select the top-1 for greedy strategies.

    \item \texttt{dwell\_prediction}: Expected duration in seconds, $\hat{d}_t$. The robot uses this to time interventions—for example, approaching 10 seconds before the predicted departure to allow for a smooth transition.

    \item \texttt{readiness}: Ordinal engagement level (A-E). High levels (A-B) indicate the visitor is deeply engaged and should not be interrupted. Low levels (D-E) suggest the visitor is ready to move on.

    \item \texttt{trajectory}: K-step rollout for planning. The robot can use this for multi-step lookahead, choosing positions that avoid blocking future viewing angles.
\end{itemize}

These outputs are structured as JSON for easy parsing by robot control systems.

\section{EXPERIMENTS}

\subsection{Metrics}

\textbf{For Next-Step Prediction:}
\begin{itemize}
    \item \textit{Hit@1:} Accuracy of top-1 prediction (does the most likely exhibit match the ground truth?)
    \item \textit{Hit@3:} Accuracy of top-3 prediction (is the ground truth among the 3 most likely?)
    \item \textit{MRR (Mean Reciprocal Rank):} $1/rank_{predicted}$, rewarding higher-ranked correct predictions
\end{itemize}

\textbf{For Dwell Time:}
\begin{itemize}
    \item \textit{MAE (Mean Absolute Error):} Average absolute error in seconds. Intuitive and directly interpretable for robot timing decisions.
\end{itemize}

\textbf{For Attention Level:}
\begin{itemize}
    \item \textit{Accuracy:} Exact match accuracy (does the predicted level match ground truth?)
    \item \textit{Cohen's $\kappa$:} Inter-rater agreement for ordinal consistency, accounting for chance agreement
\end{itemize}

\subsection{Baselines}

We compare against strong baselines across multiple categories:

\textbf{Statistical:}
\begin{itemize}
    \item \textit{Markov Chain:} First-order transition frequency model on the exhibit graph. Predicts based solely on aggregate transition probabilities, ignoring individual history.
\end{itemize}

\textbf{Deep Learning:}
\begin{itemize}
    \item \textit{LSTM:} Sequence-to-sequence model with attention mechanism (2-layer, 128 hidden units). Trained on exhibit sequences with dwell time as auxiliary output.
\end{itemize}

\textbf{Zero-Shot LLMs:}
\begin{itemize}
    \item \textit{GPT-4o, Claude 3.5 Sonnet, Gemini 2.5 Pro:} State-of-the-art commercial LLMs evaluated with task prompts but no fine-tuning. These establish upper bounds for zero-shot performance.
\end{itemize}

\textbf{Ablation Variants:}
\begin{itemize}
    \item \textit{No Memory:} Remove both short-term and long-term memory
    \item \textit{No Topology:} Remove neighbor constraints (any exhibit can be predicted)
    \item \textit{No Feature Extractor:} Remove semantic features (only exhibit IDs)
    \item \textit{No Multi-step:} Single-step prediction only (K=1)
\end{itemize}

\subsection{Main Results}

\begin{table*}[h]
\centering
\caption{Main Results on Random Split}
\begin{tabular}{llccccc}
\toprule
Split & Method & Hit@1$\uparrow$ & MRR$\uparrow$ & Hit@3$\uparrow$ & Dwell MAE$\downarrow$ & Attn Acc$\uparrow$ \\
\midrule
\multirow{5}{*}{Random}
& Markov Chain & 45.2\% & 0.62 & 68.1\% & -- & -- \\
& LSTM & 56.4\% & 0.71 & 78.2\% & 24.5s & 54.3\% \\
& GPT-4o (Zero-Shot) & 58.3\% & 0.74 & 81.2\% & 18.3s & 61.2\% \\
& Claude 3.5 (Zero-Shot) & 59.1\% & 0.75 & 82.4\% & 17.8s & 62.8\% \\
& Gemini 2.5 (Zero-Shot) & 57.6\% & 0.73 & 80.1\% & 19.2s & 59.4\% \\
& \textbf{Ours (Full)} & \textbf{68.3\%} & \textbf{0.84} & \textbf{88.4\%} & \textbf{12.1s} & \textbf{72.4\%} \\
\midrule
\multirow{2}{*}{Leave-Subj}
& Markov Chain & 38.7\% & 0.54 & 61.2\% & -- & -- \\
& LSTM & 48.2\% & 0.63 & 71.5\% & 28.3s & 48.7\% \\
& \textbf{Ours (Full)} & \textbf{61.8\%} & \textbf{0.77} & \textbf{84.1\%} & \textbf{15.4s} & \textbf{67.2\%} \\
\bottomrule
\end{tabular}
\label{tab:main_results}
\end{table*}

\textbf{Analysis:} Our method significantly outperforms statistical and deep learning baselines. The gap is particularly large for dwell time prediction (12.1s vs 17.8s MAE), suggesting that LLM-based reasoning is effective for modeling engagement. Zero-shot LLMs perform well but still lag behind our fine-tuned model, demonstrating the value of task-specific training.

Leave-subject-out results show a moderate but acceptable performance drop, indicating that the model has learned some generalizable patterns rather than memorizing individual participants.

\subsection{Ablation Study}

\begin{table}[h]
\centering
\caption{Ablation Study Results}
\begin{tabular}{lcccc}
\toprule
Variant & Next-Step Top-1$\uparrow$ & Dwell MAE$\downarrow$ & Attn Acc$\uparrow$ & Regret$\downarrow$ \\
\midrule
\textbf{Full (Ours)} & \textbf{68.3\%} & \textbf{12.1s} & \textbf{72.4\%} & \textbf{1.2} \\
-No Memory & 54.2\% & 18.3s & 61.2\% & 2.4 \\
-No Topology & 61.8\% & 15.7s & 68.1\% & 1.8 \\
-No Feature Extractor & 64.1\% & 13.9s & 70.1\% & 1.5 \\
-No Multi-step & 65.7\% & 14.2s & 71.2\% & 1.4 \\
\bottomrule
\end{tabular}
\label{tab:ablation}
\end{table}

\textbf{Analysis:}
\begin{itemize}
    \item \textit{Memory removal causes the largest drop} (14.1\% decrease in Top-1), confirming the importance of temporal patterns. This suggests that visitor behavior is highly dependent on recent history.
    \item \textit{Topology constraint removal} also hurts performance (6.5\% decrease), indicating that spatial feasibility is a strong prior.
    \item \textit{Feature extraction} contributes but is less critical than memory and topology.
    \item \textit{Multi-step prediction} provides modest but consistent gains, validating the utility of lookahead.
\end{itemize}

\subsection{Qualitative Results}

We provide qualitative examples to illustrate the system's predictions. Fig.~\ref{fig:qualitative} shows four cases:

\begin{figure}[h]
\centering
% \includegraphics[width=0.9\textwidth]{figures/qualitative.pdf}
\fbox{\parbox{0.8\textwidth}{\centering \vspace{3cm} Figure 2: Qualitative Prediction Examples \vspace{3cm}}}
\caption{Qualitative prediction examples. (a)-(c): Correct predictions showing how memory and topology guide accurate next-step prediction. (d): Failure case where visitor backtracks to previous exhibit, violating forward-flow assumptions.}
\label{fig:qualitative}
\end{figure}

\textit{Case (a):} Visitor views花卉 paintings sequentially. The model correctly predicts the next花卉 exhibit, recognizing the thematic pattern. Memory shows the visitor spent high attention (level A) on previous花卉 exhibits.

\textit{Case (b):} Visitor enters a new zone with both paintings and sculptures. Despite equal spatial proximity, the model correctly predicts a painting based on the visitor's history (previously viewed 12 paintings vs 2 sculptures).

\textit{Case (c):} Visitor lingers at an interactive display (dwell time 120s, level A). The model predicts a brief visit (level E) to the next exhibit, correctly anticipating that the visitor is fatigued and will scan remaining exhibits quickly.

\textit{Case (d):} Failure case—the visitor returns to a previous exhibit to compare details. The model predicts forward movement, missing the backtrack. This is a known limitation of our forward-flow assumption.

\section{ROBOT-FACING TOPOLOGICAL GUIDE SIMULATION}

\subsection{Simulation Setup}

To validate the robot-readiness of our learned priors, we conduct a topological guide simulation on the exhibit graph (Experiment-R).

\textbf{Environment:} The exhibit adjacency graph from OS.xls with 19 exhibits (nodes) and 37 spatial edges. This is a real exhibition space with known visitor traffic patterns.

\textbf{Human Trajectories:} Ground-truth sequences from our eye-tracking dataset, including exhibit visit order ($e_1 \rightarrow e_2 \rightarrow ... \rightarrow e_T$) and per-exhibit dwell times ($d_1, d_2, ..., d_T$).

\textbf{Robot Policies (Compared):}
\begin{itemize}
    \item \textit{Random:} Uniform random selection from neighbors—serves as a lower bound
    \item \textit{Markov-Frequency:} Select neighbor with highest historical transition frequency
    \item \textit{Shortest-Path:} Select neighbor minimizing graph distance to unvisited exhibits (requires goal inference)
    \item \textit{Ours-Prior:} Select neighbor using our model's $p(goal|context)$ distribution
\end{itemize}

\subsection{Evaluation Metrics}

\textbf{Goal Prediction Success:} Does the robot's selected goal match the human's actual next exhibit? This is directly measured by Hit@1.

\textbf{Intervention Timing Error:} Absolute difference between robot-predicted dwell and actual dwell: $|d_{pred} - d_{true}|$. Lower error means the robot can time its interventions more appropriately.

\textbf{Policy Regret:} Additional graph steps compared to optimal (ground-truth) path: $steps_{policy} - steps_{optimal}$. Regret measures how much ``wasted motion'' the robot makes due to incorrect predictions.

\subsection{Results}

\begin{table}[h]
\centering
\caption{Guide Simulation Results}
\begin{tabular}{lcccc}
\toprule
Policy & Goal Hit@1$\uparrow$ & Goal Hit@3$\uparrow$ & Timing Error$\downarrow$ & Regret$\downarrow$ \\
\midrule
Random & 15.3\% & 42.1\% & 58.4s & 4.8 \\
Markov-Frequency & 45.2\% & 68.1\% & 32.1s & 2.6 \\
Shortest-Path & 38.7\% & 62.3\% & 45.8s & 1.9 \\
\textbf{Ours-Prior} & \textbf{68.3\%} & \textbf{88.4\%} & \textbf{14.7s} & \textbf{1.2} \\
\bottomrule
\end{tabular}
\label{tab:simulation}
\end{table}

\textbf{Interpretation:}
\begin{itemize}
    \item Our prior achieves \textbf{$\sim$23\% higher goal hit rate} over the strongest baseline (Markov)
    \item \textbf{Timing error is halved} (14.7s vs 32.1s), enabling more appropriate intervention
    \item \textbf{Policy regret is lowest} (1.2 vs 2.6), indicating more human-like path following
\end{itemize}

These results demonstrate that our learned priors are not just statistically significant but robot-relevant—they can meaningfully improve guide robot decision-making.

\section{DISCUSSION}

\subsection{Why It Works}

Our approach succeeds by combining three complementary sources of information:

\textbf{Structured Topological Constraints:} The neighbor graph eliminates infeasible predictions and encodes spatial flow patterns. Even simple heuristics like ``don't jump across the exhibition space'' dramatically reduce the search space.

\textbf{Temporal Memory Patterns:} Short-term memory captures sequential dynamics (scanning patterns, re-examination), while long-term memory captures individual preferences and global exhibition structure. The 14.1\% drop when memory is removed confirms that temporal patterns are crucial.

\textbf{LLM Semantic Reasoning:} The fine-tuned LLM learns to interpret gaze through semantic coherence (related topics) and exhibit characteristics (visual salience, informational density). This enables predictions that go beyond statistical patterns—for example, recognizing that a visitor interested in landscape paintings will preferentially view other landscapes.

\subsection{Failure Modes}

Analysis of prediction errors reveals common failure cases:

\textbf{Backtracking:} Visitors sometimes revisit previous exhibits (re-examination or comparison), violating forward-flow assumptions. Our model assumes primarily forward motion through the exhibition.

\textbf{Social Interruptions:} Conversations with companions disrupt gaze patterns, causing attention to decouple from exhibit features. The model may incorrectly interpret distracted gaze as low engagement.

\textbf{Multi-Exhibit Fixations:} Visitors simultaneously viewing multiple exhibits (e.g., comparing pieces placed together) are harder to attribute. Our model assigns each gaze to a single exhibit.

\textbf{Exhibit Edge Cases:} Very short visits ($<$5s) and very long visits ($>$200s) have higher prediction variance. Short visits may be accidental glances; long visits may reflect factors not captured in our features.

\subsection{Robot Integration Outlook}

The learned priors can be integrated into robot systems in several ways:

\textbf{Goal Following:} The robot uses $p(goal|context)$ to select a following position near the predicted next exhibit. By selecting a position that doesn't block the viewing angle, the robot can maintain proximity without intrusion.

\textbf{Intervention Timing:} Dwell predictions trigger explanation delivery at appropriate engagement windows. For example, the robot might approach when predicted remaining dwell is less than 30 seconds, allowing time for the visitor to complete their viewing.

\textbf{Collision Avoidance:} Attention readiness signals guide when to pass through a visitor's viewing zone. High engagement (level A-B) means avoid; low engagement (level D-E) means safe to pass.

\textbf{Personalization:} Long-term memory can be updated per-user for adaptive service. Frequent visitors might have personalized priors based on their historical patterns.

\subsection{Limitations}

\textbf{Dataset Scale:} 8,159 segments, while substantial, may not capture full diversity of visitor behaviors. Broader data collection could reveal edge cases not present in current data.

\textbf{Venue Specificity:} Models trained on one venue may not transfer perfectly to different exhibition layouts. Cross-venue transfer learning is an area for future work.

\textbf{Annotation Noise:} Hotkey-based segmentation has inherent temporal ambiguity—different participants might press the hotkey at slightly different points in their viewing process. We mitigate this through robust training and memory-based smoothing.

\textbf{Computational Cost:} LLM inference requires $\sim$100ms per prediction, which may limit real-time applications. This can be addressed through caching (exhibit transitions repeat frequently) or distillation to smaller models.

\textbf{Single-User Assumption:} Our current approach models individual visitors in isolation. Multi-visitor scenarios (groups, families) introduce additional complexity that is not addressed.

\section{CONCLUSION}

We presented IROS\_gaze, a memory-augmented gaze reasoning framework that predicts visitor goals, dwell times, and attention levels for museum guide robots. By combining topological constraints, dual-layer memory, and LLM-based reasoning, our system achieves strong performance on next-step prediction tasks and demonstrates utility in a robot-facing simulation.

The key technical insight is that gaze, when combined with spatial and semantic context, provides a powerful prior for human intention prediction in structured environments. The dual-layer memory architecture captures both immediate sequential patterns and long-term preferences, while LLM reasoning enables flexible interpretation of complex gaze patterns.

Our robot-facing simulation demonstrates that these predictions translate to improved robot decision-making: higher goal hit rates, better timing, and lower policy regret. This validation is crucial—predictive accuracy alone does not guarantee utility for robot tasks.

Future work will focus on: (1) expanding dataset coverage to more diverse venues and visitor populations, (2) improving cross-venue generalization through domain adaptation techniques, (3) integrating the priors into real robot platforms for in-the-wild validation, and (4) extending to multi-visitor scenarios where group dynamics influence individual behavior.

The learned priors enable more socially appropriate robot navigation and intervention timing, bringing us closer to robots that can seamlessly assist visitors in complex public spaces.

\begin{thebibliography}{99}

\bibitem{shon2015people}
A. P. Shon, K. Grochow, A. Haddadi, and M. J. Matari{\'{e}}a, ``People-aware robotics: Gaze-based human intention prediction for collaborative robots,'' in \textit{Proc. IEEE Int. Conf. Robot. Autom. (ICRA)}, 2015, pp. 3520--3527.

\bibitem{fathi2012learning}
A. Fathi, Y. Li, and J. M. Rehg, ``Learning to predict daily actions using eye-gaze,'' in \textit{Proc. Eur. Conf. Comput. Vis. (ECCV)}. Springer, 2012, pp. 314--327.

\bibitem{yan2020gaze}
C. Yan, Y. Xie, F. Wang, Y. Zhu, W. Zhang, and Y. Lin, ``Gaze-based prediction of intent in supermarket shopping,'' in \textit{Proc. ACM Int. Conf. Multimodal Interact.}, 2020, pp. 305--313.

\bibitem{kitani2012activity}
K. M. Kitani, B. D. Ziebart, J. A. Bagnell, and M. Hebert, ``Activity forecasting,'' in \textit{Proc. Eur. Conf. Comput. Vis. (ECCV)}. Springer, 2012, pp. 201--214.

\bibitem{walter2024learning}
M. R. Walter, S. Hirsh, and S. Teller, ``A learning-based framework for topological map building and planning,'' in \textit{Proc. IEEE Int. Conf. Robot. Autom. (ICRA)}, 2024, pp. 11223--11230.

\bibitem{ziebart2019navigate}
B. D. Ziebart, A. L. Maas, A. K. Dey, and J. A. Bagnell, ``Navigate like a human: Intent-driven trajectory prediction,'' in \textit{Proc. IEEE/RSJ Int. Conf. Intell. Robots Syst. (IROS)}, 2019, pp. 2363--2370.

\bibitem{kucner2017semantic}
T. Kucner, M. Magnusson, and A. J. Lilienthal, ``Semantic information gain for active mapping using topological representations,'' in \textit{Proc. Eur. Conf. Mobile Robots (ECMR)}, 2017, pp. 1--8.

\bibitem{trautman2010unfreezing}
P. Trautman and A. Krause, ``Unfreezing the robot: Navigation in dense, interacting crowds,'' in \textit{Proc. IEEE/RSJ Int. Conf. Intell. Robots Syst. (IROS)}, 2010, pp. 797--803.

\bibitem{chen2017socially}
Y. F. Chen, M. Everett, M. Liu, and J. P. How, ``Socially aware motion planning with deep reinforcement learning,'' in \textit{Proc. IEEE/RSJ Int. Conf. Intell. Robots Syst. (IROS)}, 2017, pp. 1343--1350.

\bibitem{pynadath2021plan}
D. V. Pynadath and M. S. Pynadath, ``Plan recognition in multi-agent systems: A survey,'' \textit{Auton. Agents Multi-Agent Syst.}, vol. 42, no. 3, pp. 475--528, 2021.

\bibitem{liang2024spatial}
W. Liang, Y. Yue, L. Feng, and L. Feng, ``Spatial reasoning with large language models: A survey,'' \textit{ACM Comput. Surv.}, vol. 57, no. 3, pp. 1--32, 2024.

\bibitem{xia2024chain}
L. Xia, Y. Zhang, L. Shao, and H. Liu, ``Chain of thought prompting for planning with large language models,'' in \textit{Proc. AAAI Conf. Artif. Intell.}, 2024, pp. 13456--13464.

\end{thebibliography}

\end{document}
